\section{报告缺口与研究未决问题分开登记}
\label{sec:gap_map}
本节是2026年11月13日的实际复查结论。\textbf{“报告缺口”意味着现有原文已有证据而旧正文未讲清；“研究未决”意味着所读原文仍未建立该结论。}前者本轮已修订，后者必须保留；新检索未发现完整证据也不等于全领域无人解决。除下列新增Drake2009全文外，本轮对旧文是针对正文、方法与关键图页的复查，不将其重复计为新增全文核读。
\begingroup\small
\begin{longtable}{P{30mm}P{59mm}P{59mm}}
\toprule 问题与证据 & 旧报告的不足、本轮已补 & 仍存边界及判断归属\\\midrule\endfirsthead
\toprule 问题与证据 & 旧报告的不足、本轮已补 & 仍存边界及判断归属\\\midrule\endhead
双温Hall热点；Browning2016（L223）图2、3 & 补磁张力出流、离子黏性、电子超电阻及热输运因果链；不再把电子峰概括为简单欧姆热点 & 【原文】初步双温、超电阻影响强。【综合判断】峰位相容未唯一验证耗散闭合。\\
离子梯度机制；Yoon2018、2021（L272）、F743 & 补势谷/势丘、漂移层级、$\alpha$与带符号局部不稳定性；分清轨道分离和温度 & 【原文】不同导场/区域仍有其他通道。【综合判断】$\alpha>1$、轨道混沌与熵产不能直接画等号。\\
拾取与轨道混合；Drake2009Pickup、Yoon2023\allowbreak PhaseSpace & 新核读14页含附录；补二阶矩公式、功归因分歧与共同轨道积分检验 & 【原文】2009标度系数不足；2023机制份额未定。【综合判断】需同轨道预算，当前未裁决分歧。\\
电子相混；E512式(76)--(84)、Loureiro2013图2、4 & 补流体自由能到高阶矩再碰撞耗散；解释热化延迟与分离线位置 & 【原文】L244质疑无延迟与过耗散；Numata保留速度收敛限制。【综合判断】不能由相近重联率验证热化时序。\\
MRX诊断；Yoo2013图2--4、F679第II节 & 纠正电势来源；补谱线拟合、分辨率、放电筛选和温升先后；明确经典黏性解释的局限 & 【原文】热化有竞争过程，需模拟辅助。【综合判断】谱线一维拟合未直接测完整熵；与MAST黏性解释不可混用。\\
MRX分配；F674图7、Yoo2017图10 & 补控制体、损失反演、误差和大电场但小离子功的原因；纠正图10物理页码 & 【原文】中性损失取上界、热流依模型。【综合判断】实验功率比与不可逆热化比仍不同。\\
MAST温升；Tanabe2015图1--5 & 补单次时间帧与跨脉冲二维图差别、杂质诊断、EFIT时刻；串联双峰到三峰 & 【原文】缺直接磁探针、峰尖化有竞争解释。【综合判断】驻留与热输运可能掩盖局地源项变化。\\
HIT-SI与MMS；Kaptanoglu2020、TenBarge2019MMS & 保留装置几何/驱动损失及MMS各向异性侧别成败，明确是局部验证 & 【原文】电路/边界与闭合限制。【文献缺口】勘误正文、正式版新增图未核读；不能用精确比例作强结论。\\
\bottomrule\end{longtable}\endgroup

\section{从理论到观测：用什么现象区分解释}
\begin{longtable}{P{35mm}P{55mm}P{58mm}}
\toprule 物理图像 & 已有现象与原文定位 & 还需什么才能加强判定（本报告建议）\\\midrule\endfirsthead
\toprule 物理图像 & 已有现象与原文定位 & 还需什么才能加强判定（本报告建议）\\\midrule\endhead
电势先加速，后混合/碰撞 & MRX流速先升、下游谱线才明显展宽；Yoo2013图4、F679图10 & 同时看平均流、完整分布与停留时间；只拟合$T_i$无法区分冷多束和真正分布展宽。\\
出流黏性/压缩、物种交换 & L164图4的三段结构；MAST图3、4不同初始热点、晚时三峰 & 对齐几何、热边界和交换时间，测速度梯度及热流；压缩升温中可逆部分另列。\\
电子相混后弱碰撞耗散 & Loureiro2013图1、2、4的延迟和高阶谱；L244图4没有相同延迟 & 这是理论/数值对照，尚非所列装置直接观测电子Hermite级联；需分布和耗散时序验证。\\
压力张量与热流闭合 & TenBarge2019图1有物种内能；图6--10各向异性量级合理、侧别错误 & 同时匹配侧别、峰宽、热流和能量守恒，不能凭一个峰值宣布十矩已充分验证。\\
局地电功与非局地输运 & F679图11、12电功与温度峰不同；Yoo2017图8、9抗磁功大 & 联合压力功与焓/热流预算；正$\J_s\cdot\E$仅证明场向粒子输能，不证明局地不可逆热化。\\
\bottomrule\end{longtable}

上述对应也显示覆盖的真实不足：本专题对MRX/MAST的联系最具体，HIT-SI是装置整体模型比较，MMS主要依TenBarge文的事件对照；并未独立全文核读所有原始MMS观测、所有诊断方法论文及其他装置的加热系列。强导场拾取阈值、Fermi/俘获等更广的动理学加热综述仍不能由这一专题替代。它们应列为定向追读，而不是用未读摘要填进“已验证机制”。
